% !PRINTED 577
\weberpage{577}

%<<CONT
in which the summation over $\nu$ extends from 0 to $\mu$, but may
also be extended arbitrarily further, since $(\mu, \nu)$ vanishes
when $\nu$ is negative or greater than $\mu$. Accordingly we also have

\[ f(x, \mu) = -\overset{\nu}{\Sigma} (-1)^{\nu} (\mu, \nu + 1), \]

and if we therefore multiply formula (9) by $(-1)^{\nu}$ and
sum with respect to $\nu$, there results

\weq{11}{(1 - x^{\mu + 1}) f(x, \mu) = f(x, \mu + 2).}

From this it follows that $f(x, \mu)$ vanishes for every odd $\mu$, since $f(x, 1) = 0$. On the other hand, there results directly
from (10)

\[ f(x, 2) = 1 - x, \]

and therefore from (11), for every even $\mu$,

\weq{12}{f(x, \mu) = (1 - x)(1 - x^3)(1 - x^5) \ldots (1 - x^{\mu - 1}).}

For if this formula is correct for $\mu$, its correctness for
$\mu + 2$ follows from (11).

We now apply this identical transformation to our
problem, by setting $\mu = m = n - 1$ and $x = r$. Then

\[ \frac{1 - x^m}{1 - x} = \frac{1 - r^{-1}}{1 - r} = -r^{-1}, \quad \frac{1 - x^{m - 1}}{1 - x^2} = \frac{1 - r^{-2}}{1 - r^2} = -r^{-2} \cdots, \]

hence, so long as $0 \leqq \nu \leqq m$,

\[ (m, \nu) = (-1)^{\nu} r^{-(1 + 2 + \cdots + \nu)} = (-1)^{\nu} r^{-\frac{\nu(\nu + 1)}{2}}, \]

or, replacing $r$ by $r^{-2}$,

\[ (m, \nu) = (-1)^{\nu} r^{\nu^2 + \nu}. \]

Now

\[ \nu^2 + \nu \equiv \nu^2 - (n - 1) \nu \equiv \left( \nu - \frac{n - 1}{2} \right)^2 - \left( \frac{n - 1}{2} \right)^2 \ (\mathrm{mod}\ n), \]

hence

\[ (m, \nu) = (-1)^{\nu} r^{-\frac{m^2}{4}} r^{\left( \nu - \frac{m}{2} \right)^2}, \]

and it therefore becomes, by the defining equation (10),

\[ f(r^{-2}, m) = r^{-\frac{m^2}{4}} \Sigma r^{\left( \nu - \frac{m}{2} \right)^2}, \]

in which $\nu$ runs through a complete system of residues with respect to the modulus $n$.
Now simultaneously with $\nu$, $\nu - {}^1/_2 m$ also runs through such a
system of residues, so that we may also write

\[ f(r^{-2}, m) = r^{-\frac{m^2}{4}} \Sigma r^{\nu^2} \]
%CONT>>
